% Options for packages loaded elsewhere
\PassOptionsToPackage{unicode}{hyperref}
\PassOptionsToPackage{hyphens}{url}
%
\documentclass[
  10pt,
]{article}
\usepackage{amsmath,amssymb}
\usepackage{iftex}
\ifPDFTeX
  \usepackage[T1]{fontenc}
  \usepackage[utf8]{inputenc}
  \usepackage{textcomp} % provide euro and other symbols
\else % if luatex or xetex
  \usepackage{unicode-math} % this also loads fontspec
  \defaultfontfeatures{Scale=MatchLowercase}
  \defaultfontfeatures[\rmfamily]{Ligatures=TeX,Scale=1}
\fi
\usepackage{lmodern}
\ifPDFTeX\else
  % xetex/luatex font selection
\fi
% Use upquote if available, for straight quotes in verbatim environments
\IfFileExists{upquote.sty}{\usepackage{upquote}}{}
\IfFileExists{microtype.sty}{% use microtype if available
  \usepackage[]{microtype}
  \UseMicrotypeSet[protrusion]{basicmath} % disable protrusion for tt fonts
}{}
\makeatletter
\@ifundefined{KOMAClassName}{% if non-KOMA class
  \IfFileExists{parskip.sty}{%
    \usepackage{parskip}
  }{% else
    \setlength{\parindent}{0pt}
    \setlength{\parskip}{6pt plus 2pt minus 1pt}}
}{% if KOMA class
  \KOMAoptions{parskip=half}}
\makeatother
\usepackage{xcolor}
\usepackage[margin=0.75in]{geometry}
\usepackage{longtable,booktabs,array}
\usepackage{calc} % for calculating minipage widths
% Correct order of tables after \paragraph or \subparagraph
\usepackage{etoolbox}
\makeatletter
\patchcmd\longtable{\par}{\if@noskipsec\mbox{}\fi\par}{}{}
\makeatother
% Allow footnotes in longtable head/foot
\IfFileExists{footnotehyper.sty}{\usepackage{footnotehyper}}{\usepackage{footnote}}
\makesavenoteenv{longtable}
\setlength{\emergencystretch}{3em} % prevent overfull lines
\providecommand{\tightlist}{%
  \setlength{\itemsep}{0pt}\setlength{\parskip}{0pt}}
\setcounter{secnumdepth}{-\maxdimen} % remove section numbering
\ifLuaTeX
\usepackage[bidi=basic]{babel}
\else
\usepackage[bidi=default]{babel}
\fi
\babelprovide[main,import]{english}
% get rid of language-specific shorthands (see #6817):
\let\LanguageShortHands\languageshorthands
\def\languageshorthands#1{}
\ifLuaTeX
  \usepackage{selnolig}  % disable illegal ligatures
\fi
\IfFileExists{bookmark.sty}{\usepackage{bookmark}}{\usepackage{hyperref}}
\IfFileExists{xurl.sty}{\usepackage{xurl}}{} % add URL line breaks if available
\urlstyle{same}
\hypersetup{
  pdflang={en},
  hidelinks,
  pdfcreator={LaTeX via pandoc}}

\author{}
\date{}

\begin{document}

\hypertarget{the-inverted-buoyancy-origin-of-the-time-gradient-field}{%
\section{The Inverted Buoyancy Origin of the Time-Gradient
Field}\label{the-inverted-buoyancy-origin-of-the-time-gradient-field}}

Gravity as directional bias, not force

Richard Kent Gates

mail@richardkentgates.com

September 28, 2026

\hypertarget{purpose}{%
\subsection{1. Purpose}\label{purpose}}

The Time-Gradient Field Model is stated across this body of work in
terms of three postulates: a Field Postulate, a Motion Postulate, and an
Efficiency Postulate. Those postulates carry the theory. The question of
where they come from, and what the underlying picture of gravity is, has
been implicit rather than stated. This paper supplies that statement and
verifies the two properties of it that are checkable without new data.

The origin is an inverted buoyancy. An object in a medium with a density
gradient drifts toward the less dense region, because it displaces
medium and the displaced medium has a net weight. Gravity runs the
opposite way: matter moves toward \emph{slower} time, not faster. Read
the proper-time rate as the density of a medium, and the sign inverts.
This is the intuition that motivates the Motion Postulate and fixes its
direction. It is not the mechanism, and this paper does not treat it as
one.

\hypertarget{the-motion-postulate-as-a-gradient-rule}{%
\subsection{2. The Motion Postulate as a gradient
rule}\label{the-motion-postulate-as-a-gradient-rule}}

The formulation that the papers actually use is a gradient-following
rule, stated in \emph{Mathematical Verification of the Time-Gradient
Field Model} as follows:

"This is not a force but a directional bias from the gradient ∂Φ/∂x."

An object accelerates along the local gradient of the time field,
evaluated at the object\textquotesingle s own location. There is no
force term, and no integral over the body\textquotesingle s volume. This
distinction is not cosmetic. A volume integral over an extended body
would raise the question of how the body\textquotesingle s internal mass
distribution affects the result, and that question does not arise here,
because the response is local.

For a spherical source of mass \textbackslash(M\textbackslash), the
proper-time rate is

\textbackslash{[}\textbackslash frac\{d\textbackslash tau\}\{dt\} =
\textbackslash sqrt\{1 -
\textbackslash frac\{2GM\}\{rc\^{}2\}\}\textbackslash{]}

so the gradient at radius \textbackslash(r\textbackslash) is

\textbackslash{[}\textbackslash frac\{\textbackslash partial
(d\textbackslash tau/dt)\}\{\textbackslash partial r\}
\textbackslash approx \textbackslash frac\{GM\}\{r\^{}2
c\^{}2\}\textbackslash{]}

and the acceleration, restoring the factor
\textbackslash(c\^{}2\textbackslash) that converts a rate of change of a
time ratio into an acceleration, is

\textbackslash{[}a = c\^{}2 \textbackslash frac\{\textbackslash partial
(d\textbackslash tau/dt)\}\{\textbackslash partial r\} =
\textbackslash frac\{GM\}\{r\^{}2\}\textbackslash{]}

This is Newtonian acceleration, recovered identically rather than
approximately, because
\textbackslash(d\textbackslash tau/dt\textbackslash) already contains
\textbackslash(GM/r\^{}2\textbackslash). The analogy does not need to
approximate the inverse-square law; the law is a restatement of the
gradient.

\hypertarget{verification}{%
\subsection{3. Verification}\label{verification}}

Both properties below are computed in \texttt{buoyancy\_origin.py},
which writes \texttt{buoyancy\_origin\_results.txt}. Neither requires
new data: both are algebraic and geometric consequences of the
prescription.

\hypertarget{the-gradient-reproduces-newtonian-acceleration}{%
\subsubsection{3.1 The gradient reproduces Newtonian
acceleration}\label{the-gradient-reproduces-newtonian-acceleration}}

Evaluating the gradient numerically by central difference and comparing
against \textbackslash(GM/r\^{}2\textbackslash), at the
Earth\textquotesingle s surface:

\begin{longtable}[]{@{}ll@{}}
\toprule\noalign{}
Quantity & Value \\
\midrule\noalign{}
\endhead
\bottomrule\noalign{}
\endlastfoot
\textbackslash(d\textbackslash tau/dt\textbackslash) at \textbackslash(r
= 6.371\textbackslash times10\^{}6\textbackslash) m &
\textbackslash(9.99999999304\textbackslash times10\^{}\{-1\}\textbackslash) \\
Weak-field parameter \textbackslash(GM/(rc\^{}2)\textbackslash) &
\textbackslash(6.960\textbackslash times10\^{}\{-10\}\textbackslash) \\
Numerical gradient &
\textbackslash(1.092515317\textbackslash times10\^{}\{-16\}\textbackslash)
s\textbackslash(\^{}\{-1\}\textbackslash) \\
Analytic gradient &
\textbackslash(1.092515316\textbackslash times10\^{}\{-16\}\textbackslash)
s\textbackslash(\^{}\{-1\}\textbackslash) \\
Acceleration from gradient & \textbackslash(9.8195320403\textbackslash)
m s\textbackslash(\^{}\{-2\}\textbackslash) \\
Analytic \textbackslash(GM/r\^{}2\textbackslash) &
\textbackslash(9.8195320328\textbackslash) m
s\textbackslash(\^{}\{-2\}\textbackslash) \\
Relative error &
\textbackslash(7.580\textbackslash times10\^{}\{-10\}\textbackslash) \\
Result & PASS \\
\end{longtable}

The weak-field linearization of
\textbackslash(d\textbackslash tau/dt\textbackslash) introduces no
measurable error at planetary radii: the linearized and exact gradients
differ by
\textbackslash(6.960\textbackslash times10\^{}\{-10\}\textbackslash) in
relative terms, the same order as the parameter itself.

Implementation note. The quantity \textbackslash(1 -
d\textbackslash tau/dt\textbackslash) must be evaluated as
\textbackslash(x/(1 + \textbackslash sqrt\{1-x\})\textbackslash) with
\textbackslash(x = 2GM/(rc\^{}2)\textbackslash). The naive form
\textbackslash(1 - \textbackslash sqrt\{1-x\}\textbackslash) subtracts
two numbers near unity whose difference is of order
\textbackslash(10\^{}\{-9\}\textbackslash), which discards most of the
available significant figures and returns zero under double-precision
arithmetic.

\hypertarget{the-coupling-carries-no-compositional-information}{%
\subsubsection{3.2 The coupling carries no compositional
information}\label{the-coupling-carries-no-compositional-information}}

The time field is a property of the source mass. It contains no
information about the material of the body responding to it. An object
therefore follows the same gradient regardless of what it is made of,
and the acceleration depends on the medium, not on the body. Evaluating
the gradient at a single radius for bodies spanning twenty orders of
magnitude in density:

\begin{longtable}[]{@{}llll@{}}
\toprule\noalign{}
Body & Density (kg m\textbackslash(\^{}\{-3\}\textbackslash)) & Mean
\textbackslash(Z\textbackslash) &
\textbackslash(d\textbackslash tau/dt\textbackslash) gradient
(s\textbackslash(\^{}\{-1\}\textbackslash)) \\
\midrule\noalign{}
\endhead
\bottomrule\noalign{}
\endlastfoot
Hydrogen (H\textsubscript{2}) &
\textbackslash(8.99\textbackslash times10\^{}\{-8\}\textbackslash) &
0.67 &
\textbackslash(2.7312882907\textbackslash times10\^{}\{-17\}\textbackslash) \\
Water (H\textsubscript{2}O) &
\textbackslash(1.00\textbackslash times10\^{}\{3\}\textbackslash) & 3.33
&
\textbackslash(2.7312882907\textbackslash times10\^{}\{-17\}\textbackslash) \\
Silicate rock &
\textbackslash(2.50\textbackslash times10\^{}\{3\}\textbackslash) &
10.00 &
\textbackslash(2.7312882907\textbackslash times10\^{}\{-17\}\textbackslash) \\
Iron core &
\textbackslash(7.87\textbackslash times10\^{}\{3\}\textbackslash) &
26.00 &
\textbackslash(2.7312882907\textbackslash times10\^{}\{-17\}\textbackslash) \\
Lead & \textbackslash(1.34\textbackslash times10\^{}\{4\}\textbackslash)
& 82.00 &
\textbackslash(2.7312882907\textbackslash times10\^{}\{-17\}\textbackslash) \\
Degenerate matter &
\textbackslash(1.00\textbackslash times10\^{}\{12\}\textbackslash) &
2.00 &
\textbackslash(2.7312882907\textbackslash times10\^{}\{-17\}\textbackslash) \\
\multicolumn{2}{@{}l}{%
Spread across all compositions} &
\textbackslash(0.000\textbackslash times10\^{}\{0\}\textbackslash) & \\
\multicolumn{2}{@{}l}{%
Result} & PASS & \\
\end{longtable}

Mean \textbackslash(Z\textbackslash) is the atomic number per
constituent nucleus, averaged by number: H\textsubscript{2} gives
\textbackslash(2/3 = 0.67\textbackslash), H\textsubscript{2}O gives
\textbackslash((2\textbackslash times1 + 8)/3 = 3.33\textbackslash), and
SiO\textsubscript{2} gives \textbackslash((14 +
2\textbackslash times8)/3 = 10.00\textbackslash). The final column is
the gradient of the proper-time rate, in
s\textbackslash(\^{}\{-1\}\textbackslash); the acceleration is this
value multiplied by \textbackslash(c\^{}\{2\}\textbackslash).

The spread is exactly zero. The equivalence principle holds for the
origin on its own terms, with no appeal to the equivalence principle as
an input.

\hypertarget{consequence-for-the-sector-coupling}{%
\subsection{4. Consequence for the sector
coupling}\label{consequence-for-the-sector-coupling}}

The gravity-sector coupling in the synthesis is written

\textbackslash{[}c\_g = \textbackslash sum\_i f\_i \textbackslash left(
d\_\textbackslash alpha + \textbackslash beta\_i d\_\textbackslash mu
\textbackslash right)\textbackslash{]}

where \textbackslash(f\_i\textbackslash) are material composition
fractions and \textbackslash(\textbackslash beta\_i\textbackslash) are
composition-dependent sensitivity coefficients. Section 3.2 establishes
that the buoyancy origin alone predicts \textbackslash(c\_g =
0\textbackslash) for the composition-dependent part, because the medium
carries no compositional information. Any observed value of
\textbackslash(c\_g\textbackslash) that varies with composition is
therefore a \emph{measured deviation from a principled baseline}, not a
parameter introduced to fit data.

This is a stronger claim than a fit, and it is testable in a way a fit
is not. A composition-independent result would be consistent with the
origin. A composition-dependent result would require the deviation to be
attributed to something the origin does not contain, and that something
would be new physics with a specific place to attach.

\hypertarget{what-this-paper-does-not-establish}{%
\subsection{5. What this paper does not
establish}\label{what-this-paper-does-not-establish}}

\textbf{The entropy arrow.} The gradient picture supplies a direction
from a gradient. It does not, by itself, supply a time asymmetry: the
field equation is time-reversible, and reversibility is what allows the
gradient to have a sign under time reversal. In the papers the entropy
arrow is asserted as the origin of that sign. Establishing it is
separate work, not verification of an existing claim, and it is the most
consequential open item in the framework.

\textbf{The non-gravitational sectors.} The buoyancy origin yields
gravity. It does not by itself yield the decay-rate coupling, the
fine-structure variation, or the mass-ratio variation. The unified
correction law \textbackslash(X\_\{\textbackslash text\{eff\}\} = X\_0(1
+ c\_X G)\textbackslash) is a postulate laid across the sectors, not a
consequence of the gradient picture. Whether one field can serve all
sectors is the central open question, and this paper does not address
it.

\textbf{Datasets.} No anomaly measurement is used, fitted, or checked
here. The evidence base is documented in the companion papers on the
time-gradient field model and its statistical verification.

\hypertarget{conclusion}{%
\subsection{6. Conclusion}\label{conclusion}}

Gravity in this framework is a directional bias along the local gradient
of proper-time rate, not a force. Inverted buoyancy supplies the
intuition for its sign: ordinary buoyancy carries a body toward less
dense medium, gravity carries matter toward slower time, so the sign
inverts while the mechanism of a density gradient carrying a body with
it does not.

Two properties of this origin are established. The gradient of the
proper-time rate reproduces Newtonian acceleration to a relative error
below \textbackslash(10\^{}\{-9\}\textbackslash), with no measurable
contribution from linearization at planetary radii. The second follows
from the structure of the model rather than from measurement: the
coupling carries no compositional input, so the equivalence principle
holds exactly and composition dependence in
\textbackslash(c\_g\textbackslash) becomes a measurable deviation rather
than a fitted parameter.

\href{https://richardkentgates.com/}{← richardkentgates.com}

\href{https://richardkentgates.com/}{← richardkentgates.com}

\hypertarget{ai-research-collaboration-disclosure}{%
\subsection{AI Research Collaboration
Disclosure}\label{ai-research-collaboration-disclosure}}

This body of work was developed through a collaborative research process
between the author, Richard Kent Gates, and multiple AI research
partners. The author provides full transparency on this process.

\textbf{Author\textquotesingle s Role (Richard Kent Gates):}

\begin{itemize}
\tightlist
\item
  Conceived all core theories, conceptual frameworks, hypotheses, and
  research direction
\item
  Defined the Time-Gradient Field Model, the unified scalar field G(t),
  the distance → time → information → G(t) framework, and all
  theoretical architecture
\item
  Provided physical intuition, biological reasoning, and domain
  expertise that guided every theoretical decision
\item
  Reviewed, validated, and approved all mathematical formalisms before
  inclusion
\item
  Made all final judgments on theoretical coherence and physical
  interpretation
\end{itemize}

\textbf{AI Research Partners:}

\begin{itemize}
\tightlist
\item
  \textbf{Mimo V2.5 (opencode platform)} --- Primary mathematical
  formalization partner. Assisted in translating conceptual physics into
  mathematical notation, generating numerical proof scripts, compiling
  LaTeX documents, and building the website infrastructure.
\item
  \textbf{Microsoft Copilot (browser-based)} --- Assisted in literature
  review, dataset gathering, cross-referencing empirical results (DESI,
  BNL, PTB, MICROSCOPE), and verifying citation accuracy.
\item
  \textbf{Google Gemini (browser-based)} --- Assisted in conceptual
  exploration, thought experimentation, and early-stage hypothesis
  refinement.
\item
  \textbf{Space Bunny (OpenCode)} --- Verification and provenance
  auditing. Read every paper and proof script in the repository and
  traced each reported figure back to the script and line that produces
  it. Independently recomputed the Fisher information matrices, the
  Cramér--Rao bounds, and the parameter errors from the models as
  stated. Resolved the origin of the Fisher information values quoted in
  the unified and Theory-of-Everything papers, reconstructed the
  two-parameter Fisher matrix in the statistical verification paper from
  its sampling grid, and identified a mislabeled citation and a
  statistical criterion stated more narrowly than the reported results
  supported. No theoretical claim, numerical result, or interpretive
  judgment in this body of work originated with this tool.
\end{itemize}

\textbf{Nature of AI Involvement:}

The AI tools functioned as research assistants --- analogous to graduate
students or technical collaborators who help formalize, compute, and
organize ideas that originate from the principal investigator. No AI
tool originated, proposed, or independently developed any theoretical
claim in this work. All physical insights, theoretical innovations, and
interpretive judgments are the author\textquotesingle s own.

\end{document}
